\subsection{One integer for all coefficient primes in rank three}\label{rk3:sec}

Fix an imaginary quadratic field \(K\), a level \(N\ge1\), a twist
\(e\ge1\), a prime \(p\nmid30N|\operatorname{disc}K|\), and a place
\(w\mid p\). Write \(q=|\mathcal O_K/w|=p^h\), \(h=1\) or \(2\).
In this subsection \(q\) is a residue norm, independent of the notation
for the first auxiliary prime elsewhere. For \(H\le(\mathbf Z/N)^\times\),
let \(J\) be the complex characters modulo \(N\) with kernel exactly \(H\).
If the quotient is cyclic then \(|J|=\varphi([(\mathbf Z/N)^\times:H])\);
otherwise \(J\) is empty. This follows by choosing the image of a generator
under a faithful embedding in \(\mathbf C^\times\). Put \(z_j=j(q)^{-1}\).
The multiset of the \(z_j\) is invariant under every automorphism of
\(\overline{\mathbf Q}\).

The two spectral lists used at \(w\) are the irreducible rank-three,
polarized automorphic representations with normalized weight sequence
\((e,0,-e)\), and the sums \(W|_{G_K}\oplus1\) from a rank-two newform
of weight \(2e+1\), with the normalization giving weights \((e,-e)\).
The conductor bound for the second list involves \(N|\operatorname{disc}K|\).
The following lemma isolates the exact use of classification by finite
lists. The same trace can belong to several distinct entries.

\begin{proposition}[Counting the classical spectral lists]\label{rk3:count}
Fix \(F=\mathbf Q\) or an imaginary quadratic field and \(n\ge2\).
Let \(K_F(M)\) be the kernel of reduction modulo \(M\) in
\(\mathrm{GL}_n(\widehat{\mathcal O}_F)\).
Consider regular algebraic cuspidal representations of
\(\mathrm{GL}_n(\mathbb A_F)\) with \(K_F(M)\)-fixed vectors and
one prescribed infinitesimal character at each embedding. Suppose the
pairwise differences of its Harish--Chandra parameters are at most \(H\)
in absolute value. The number of isomorphism classes is at most
\[
 C_{F,n}(3M)^{n^2[F:\mathbf Q]}
       (1+H)^{[F:\mathbf Q]n(n-1)/2}.
\]
On the split real centre take the prescribed character and the inverse
character on the coefficient system; the product descends to the quotient.
\end{proposition}
\begin{proof}
The regular algebraic cuspidal condition gives nonzero relative Lie
algebra cohomology with the irreducible coefficient module of the
prescribed infinitesimal character, using the dual if necessary.
Take the connected maximal compact at infinity, together with the
split central factor. Distinct cuspidal representations give distinct
summands in cuspidal cohomology. This embeds in ordinary cohomology by
the cuspidal summand in Franke's theorem. At finite level take compact-open
invariants, an exact functor over \(\mathbf C\) by averaging. Since
\(K_F(3M)\subseteq K_F(M)\), every representation counted contributes
at least one dimension to the sum over all cohomological degrees at
level \(3M\).
Grobner--Raghuram,
Thms.~6.3 and 7.6, Cor.~7.7 and Prop.~7.16~\cite{grobnerRaghuram}.

At level three there are finitely many connected components. Indeed the
determinants give finitely many ray classes and strong approximation
for \(\mathrm{SL}_n\) gives the component description, up to the
finite group of real signs. The stabilizers are torsion-free: in each
adelic lattice their action on its underlying integral lattice is
congruent to the identity modulo three. The elementary binomial
argument in Theorem~\ref{newcoh:lemmaK} applies. The Borel--Serre
compactifications of these components have the homotopy type of the
interior and are compact manifolds with corners. Fix finite
triangulations with \(C_{F,n}\) simplices in total. The finite cover
at \(3M\) has at most \([K_F(3):K_F(3M)]\) lifts per simplex, at most
\((3M)^{n^2[F:\mathbf Q]}\) by counting all residue matrices modulo
\(3M\). A cellular cochain group with local coefficients has one
copy of the fibre per cell. Thus the sum of the Betti numbers is
bounded by the number of lifted simplices times the fibre dimension.
Borel--Serre,
\emph{Corners and arithmetic groups}~\cite{borelSerre}.

At one embedding write the ordered infinitesimal character as
\(\mu+\rho\), with \(\mu\) the highest weight. Weyl's formula for
the coefficient dimension multiplies its pairwise positive differences
divided by \(j-i\) over \(i<j\). There are \(n(n-1)/2\) factors,
each at most \(1+H\); changing a determinant twist leaves the
differences fixed. Taking the product over all embeddings proves the
bound. For the present rank-three triples the infinitesimal character
is \((e,0,-e)\). Subtracting \(\rho=(1,0,-1)\) gives dimension \(e^3\)
at each of the two embeddings, hence \(e^6\) in all. In rank two
\((e+1/2,1/2-e)-\rho=(e,1-e)\) has dimension \(2e\).
\end{proof}

\begin{corollary}[Discriminant and level in the rank-three count]\label{rk3:volume}
Let \(D=|\operatorname{disc}K|\) for an imaginary quadratic field \(K\).
Let \(S\) consist of the two classical lists of
Proposition~\ref{rk3:count}, at principal levels \(N\) over \(K\)
and \(ND\) over \(\mathbf Q\), respectively, with one entry per
system. For integers \(N,e\ge1\),
\[
 |S|\ll h_K D^4 N^{18}e^6+(ND)^4e
       \ll D^5N^{18}e^6 .
\]
The constants are absolute. Consequently the integer of
Lemma~\ref{rk3:integer}, under its trace hypotheses, satisfies
\[
 \log|C_p|\ll |J|D^5N^{18}e^7\log p .
\]
\end{corollary}
\begin{proof}
Put \(L=3N\) and \(X=\mathrm{SL}_3(\mathbf C)/\mathrm{SU}(3)\)
with a fixed invariant Riemannian metric. As in
Proposition~\ref{rk3:count}, pass to principal level \(L\) and
divide out the real split centre. At infinity
\(\mathrm{GL}_3(\mathbf C)/(\mathbf R_{>0}\mathrm U(3))\simeq X\).
Determinant and strong approximation for \(\mathrm{SL}_3\) identify
the connected components with the finite ray class group modulo
\(L\mathcal O_K\). In fact the determinant image of the principal
compact open is the subgroup of finite integral idelic units
congruent to one at that modulus, as is seen on diagonal matrices.
The map to the ordinary class group has kernel a quotient of
\((\mathcal O_K/L\mathcal O_K)^\times\). There are thus at
most \(h_K L^2\) components.

For a finite adelic representative \(g\) of a component put
\(\Lambda_v=g_v\mathcal O_{K_v}^3\). A member of its arithmetic
stabilizer preserves these lattices and acts trivially modulo \(L\).
Its determinant is a unit \(\epsilon\in\mathcal O_K^\times\) with
\(\epsilon\equiv1\bmod L\). If \(\epsilon\ne1\), then
\(0<|\operatorname{N}_{K/\mathbf Q}(\epsilon-1)|\le4\), whereas
divisibility by \(3\) makes that norm divisible by \(9\).
Hence the stabilizer is the level-\(L\) subgroup of
\[
 \Gamma_g^1=\mathrm{SL}_3(K)\cap
          \prod_{v<\infty}\operatorname{Stab}_{\mathrm{SL}_3(K_v)}
                                     (\Lambda_v).
\]
Its index in \(\Gamma_g^1\) is at most \(L^{16}\). Indeed, in a
basis of each \(\Lambda_v\) the image of reduction at
\(\mathfrak p_v^b\Vert L\mathcal O_K\) has size at most
\[
 |\mathrm{SL}_3(\mathcal O_{K_v}/\mathfrak p_v^b)|
 =q_v^{8(b-1)}|\mathrm{SL}_3(\mathbf F_{q_v})|\le q_v^{8b}.
\]
Reduction is surjective, since a lift with unit determinant can be
corrected by scaling one row. Each successive kernel is the additive
space of trace-zero matrices over the residue field. Multiplication
over \(v\) gives \(\operatorname{N}(L\mathcal O_K)^8=L^{16}\).

Prasad's covolume formula applies to the split simply connected
\(K\)-group \(\mathrm{SL}_3\). For every \(g\) each local stabilizer
is hyperspecial, with integral model \(\mathrm{SL}(\Lambda_v)\)
and reductive special fibre \(\mathrm{SL}_3\). They form a coherent
collection, equal to the standard one almost everywhere. Strong
approximation at the noncompact complex place gives the principal
arithmetic group appearing in the formula. Its Haar covolume in
\(\mathrm{SL}_3(\mathbf C)\) is
\[
 c_\infty D^{8/2}\prod_{v<\infty}
   \frac{q_v^8}{|\mathrm{SL}_3(\mathbf F_{q_v})|}
 =c_\infty D^4\zeta_K(2)\zeta_K(3).
\]
Here the Tamagawa number is one and the archimedean constant is
fixed for the single complex place and the chosen Haar measure.
In particular there is no dependence on the ideal class of \(g\).
The local zeta factors give \(\zeta_K(j)\le\zeta(j)^2\) for \(j>1\).
Prasad,
\emph{Volumes of \(S\)-arithmetic quotients of semi-simple groups},
\S\S3.3, 3.7 and 3.11~\cite{prasadVolume}.

The level-\(L\) subgroup acts faithfully and without torsion, by
its action on the underlying integral lattice modulo three. Its
quotient of \(X\) is noncompact since \(\mathrm{SL}_3\) is \(K\)-split
of positive rank. Its Riemannian volume is at most a fixed constant
times \(D^4L^{16}\), by the covolume and index bounds. Gelander's
theorem gives a homotopy-equivalent simplicial complex with at most
a fixed multiple of this volume in vertices and uniformly bounded
vertex degree. A vertex of degree \(b\) belongs to at most \(2^b\)
simplices, so the same linear bound holds for all simplices.
Local-coefficient cochains contribute at most \(e^6\) dimensions per
simplex. Summing over the \(h_K L^2\) possible components and
using cuspidal cohomology as in Proposition~\ref{rk3:count} proves
the first term. The second follows from the same proposition over
the fixed field \(\mathbf Q\) in rank two with \(M=ND\) and
coefficient dimension \(2e\).
Gelander,
\emph{Homotopy type and volume of locally symmetric manifolds},
Theorem 1.5(1)~\cite{gelanderVolume}.

Minkowski supplies each ideal class with an integral ideal of norm
at most \(\sqrt D\). There are at most \(\tau(m)\) ideals of norm
\(m\), by the quadratic prime splitting possibilities. Consequently
\(h_K\le\sqrt D(1+\log\sqrt D)\le D\).
Finally \(q\le p^2\), so the per-factor logarithmic cost
\(2e\log q+1\) of Lemma~\ref{rk3:integer} is \(O(e\log p)\).
\end{proof}


\begin{lemma}[Size and denominators of the classical traces]\label{rk3:purity}
For the conjugate self-dual rank-three cuspidal systems just described,
take principal level \(N\) and the weights for the inverse cyclotomic
twist \(m=(e-1,-1,-e-1)\) in the convention that the labelled jumps
are \(-m_i\). For the rank-two summands take principal level
\(N|\operatorname{disc}K|\) and weights \((e,-e)\). Assume the usual
classical polarized cuspidal realizations with these levels and weights.
The resulting traces at inverse arithmetic Frobenius, indexed with
multiplicity by the systems, are algebraic, stable as a list under
conjugation when all such systems are taken, and satisfy
\(|t_s|\le3\), \(q^et_s\) an algebraic integer.
\end{lemma}
\begin{proof}
Use geometric Frobenius, the inverse of arithmetic Frobenius. In rank
three the spherical parameters \(\alpha_i\) in unitary normalization
have absolute value one. Indeed the complex representation is cuspidal,
cohomological and conjugate self-dual over a CM field, so Caraiani's
temperedness theorem applies. By the normalization of
Proposition~\ref{newgal:classical}, the geometric eigenvalues of the
inverse cyclotomic twist are \(b_i=q\alpha_i\). The cyclotomic
character at geometric Frobenius is \(q^{-1}\), so restoring the
original representation gives \(t_s=q^{-1}\sum_i b_i\).
The cohomological conjugation theorem permutes the systems: all
embeddings have the same weight triple, the level is defined
arithmetically, and conjugation commutes with the contragredient and
field-automorphism relations on the finite parts. Strong multiplicity
one identifies the conjugate cuspidal representation by its finite
part. The spherical parameters in algebraic normalization are algebraic
and conjugate accordingly. Here the factor \(q\) is rational.
Caraiani,
\emph{Local-global compatibility and the action of monodromy on
nearby cycles}, Thm.~1.2~\cite{caraianiNearby};
Grobner--Raghuram,
Thm.~7.17 and Remark~7.18~\cite{grobnerRaghuram}.

At a coefficient prime \(v\nmid p\), a stable integral lattice for
the corresponding member of the compatible system shows that both the
eigenvalues of Frobenius and their inverses are integral: its matrix
and its inverse preserve the lattice. At the coefficient prime \(p\),
the local representation is crystalline since \(p\) divides neither
the level nor the discriminant. Let \(q=p^h\). The normalized Newton
slopes of its filtered isocrystal are \(v_p(b_i)/h\), with
\(v_p(p)=1\). They lie between the smallest and largest labelled
filtration jumps. To see the lower bound, take the isoclinic subobject
of minimum slope with its induced filtration, after extending
coefficients as necessary. Weak admissibility gives \(t_N\ge t_H\)
there, and each jump is at least the minimum of the ambient jumps.
Divide by its rank, with the usual average over embeddings. Apply the
same argument to the dual for the upper bound. In rank three the
endpoints are \(1-e,1+e\). Hence \(q^{e-1}b_i\) is \(p\)-integral.
The argument applies to every embedding of the coefficient field.
Integrality at every finite place gives the assertion for \(q^et_s\).

For a normalized rank-two member \(W\) over \(\mathbf Q\) the jumps
are \(\{-e,e\}\). The purity weight in the polarized cuspidal theorem
is zero, since the reflected pair of jumps about \(w/2\) must be that
same pair. Its two geometric Frobenius eigenvalues \(b_1,b_2\) at a
good \(p\) therefore have absolute value one under every embedding.
The \(p\)-adic slope calculation gives \(v_p(b_i)\ge-e\). Restriction
to \(K_w\) raises Frobenius to the \(h\)th power modulo inertia, which
acts trivially here. Thus \(t_s=1+b_1^h+b_2^h\), and integrality at
every finite place follows as above. Equivalently, if the newform is
in motivic weight \(2e\), its eigenvalues are divided by \(p^e\);
this is an integral power and commutes with conjugation. This proves
the assertions in both ranks. For the comparison between spherical
parameters, labelled jumps and crystalline Frobenius see
Barnet-Lamb--Gee--Geraghty--Taylor,
Thm.~2.1.1~\cite{blggtPotential}.
\end{proof}

\begin{lemma}[The integral trace product]\label{rk3:integer}
Suppose the finite list \((t_s)_{s\in S}\) of algebraic numbers is
Galois invariant with multiplicities, each \(|t_s|\le3\), and each
\(q^et_s\) is an algebraic integer. Set \(Q=q^e\) and
\[
 t^0_j=Q^{-1}z_j+1+Qz_j^{-1},\qquad
 C_p=\prod_{j\in J}\prod_{s\in S}Q(t_s-t^0_j).
\]
Then \(C_p\) is a nonzero rational integer with
\[
 \log|C_p|\le |J|\,|S|\,(2e\log q+1).
\]
For any \(\ell\ne p\), fix an embedding of the numbers involved into
\(\overline{\mathbf Q}_\ell\), with \(v_\ell(\ell)=1\). If at every
\(j\) a sublist of indices \(S_j\subseteq S\) has been selected, then
\[
      \sum_{j\in J}\sum_{s\in S_j}v_\ell(t_s-t^0_j)
       \ \le\ v_\ell(C_p).
\]
\end{lemma}
\begin{proof}
In a factor \(Qt_s-z_j-Q-Q^2z_j^{-1}\), all terms are algebraic
integers. The last has absolute value \(Q^2\), the sum of the others
has magnitude at most \(4Q+1<Q^2\) since \(Q\ge7\), and the magnitude
of the factor is at most \(2Q^2\). The action of any Galois automorphism
permutes the full product. A rational algebraic integer is an integer.
Sum \(\log(2Q^2)\le1+2e\log q\) over the factors, with the empty product
interpreted as one. At \(\ell\ne p\), \(Q\) is a unit. Every omitted
factor has nonnegative valuation by integrality. Removing those factors
proves the inequality. The embedding may vary with \(\ell\); the integer
in the bound remains fixed.
\end{proof}

Here are the prime and character conditions at the rank-three
application. The modulus \(N\) is even, \(e\ge1\),
\(\psi(-1)=(-1)^e\), and \(\ell\) is odd with
\(\ell\nmid\operatorname{ord}(\psi)|\operatorname{disc}K|h_K\) and
\(\ell-1\nmid 2e\). The \(\ell\)-local component of \(\psi\) has
square one and \((N,|\operatorname{disc}K|)=1\). If \(\psi^{(\ell)}\)
is obtained by deleting that local component, the quadratic field \(K\)
must not be contained in the compositum of \(\mathbf Q(\mu_\ell)\)
and the field cut out by \(\psi^{(\ell)}\). These conditions have the
same truth value for every faithful character of the cyclic image:
raising its values to a power coprime to the image order preserves
the kernels, the order of every component, and the sign at \(-1\).
\src{src/Development/AutomorphicArithmetic/AutomorphicVarieties/HeightAndVanishing/Counts.lean}{9}{Prime and character hypotheses}.

\begin{lemma}[Independence of the integral coefficient ring]\label{rk3:coeff}
Let \(L_0\) be the extension of \(\mathbf Q_\ell\) generated by the
values of the character on units, and let \(\mathcal O_0\) be its ring of integers.
For a finite extension \(L/L_0\), put \(A_L=(L/\mathcal O_L)
(\psi^{-1}\chi_\ell^{-e})\). The strict groups at level \(\ell N\)
satisfy
\[
 S_{\ell N}(\mathbf Q,A_L)\simeq
 S_{\ell N}(\mathbf Q,A_{L_0})\otimes_{\mathcal O_0}\mathcal O_L .
\]
Finiteness and length divided by the absolute ramification index
are independent of \(L\).
\end{lemma}
\begin{proof}
The character values on units have finite order prime to \(\ell\), so lie in
an unramified extension by Hensel lifting the roots of \(X^m-1\)
over a residue field containing \(\mu_m\), \(m=\operatorname{ord}\psi\).
The remaining Tate character takes values in \(\mathbf Z_\ell^\times\).
The finite free extension of valuation rings is flat, giving
\((L_0/\mathcal O_0)\otimes\mathcal O_L=L/\mathcal O_L\).
In a free \(\mathcal O_0\)-basis of \(\mathcal O_L\), each group of
continuous cochains becomes the finite direct sum of its counterparts
over \(\mathcal O_0\), with coordinatewise coboundaries and restrictions.
The common kernel of the restriction maps is obtained coordinatewise,
even for the conditions over all inertia subgroups. This proves the
displayed isomorphism. Finite length is preserved and reflected since
the extension is finite free and faithfully flat. Each elementary factor
\(\mathcal O_0/(\pi_0^k)\) gives length \(k\,e(L/L_0)\) after extension.
Dividing by \(e(L/\mathbf Q_\ell)=e(L/L_0)e(L_0/\mathbf Q_\ell)\)
proves the length assertion.
\end{proof}

To see precisely what must be supplied by the deformation argument, let
\(\mathcal O\) be a coefficient discrete valuation ring of ramification
index \(a\) over \(\mathbf Z_\ell\) and write \(\lambda_j\) for the
finite length over \(\mathcal O\) of the relevant strict Selmer group,
divided by \(a\). This normalization makes length and \(v_\ell\)
compatible under extension of coefficients. In the rank-three case the
needed assertion gives, for each \(j\), a finite set of pairwise
nonisomorphic admissible representations \(\rho\) with
\[
       \lambda_j\le
          2\sum_\rho v_\ell\bigl(\operatorname{tr}
                       \rho(\mathrm{Frob}_w^{-1})-t^0_j\bigr).
\]
It includes finiteness of the Selmer group. The trace-list theorem must
classify these representations by distinct indices within each such set,
identify their normalized inverse-Frobenius traces with the \(t_s\),
and give \(|S|\le C_K(N(e+1))^{A_K}\) with constants independent of
\(p,\ell,e,N,H\). When one character with kernel \(H\) exists over
\(\mathcal O\), every other embedding of the cyclic image is a coprime
power of that character, so the orbit used in the length sum is exactly
the complex list \(J\). The lemma now gives
\(\sum_j\lambda_j\le2v_\ell(C_p)\). It explains why one cannot replace
the spectral list by an unlabelled set of different trace values.

Choose another allowable prime \(p'\ne p\). For any finite set of
allowable coefficient primes, sum the first estimate away from \(p\),
and use \(p'\) only at \(p\). Prime factorization gives
\[
 \sum_{\ell}\log\ell\sum_j\lambda_j
       \le 2\log|C_p|+2\log|C_{p'}|
       \ll \varphi([(\mathbf Z/N)^\times:H])\,
          |\operatorname{disc}K|^5 N^{18}e^7\log(pp').
\]
This bound for finite subsums also bounds the full nonnegative sum
whenever the stated local hypotheses hold throughout it. For the
classical lists Corollary~\ref{rk3:volume} and Lemma~\ref{rk3:purity}
give the displayed exponents and the integral factors. The implied
constant is absolute. For each fixed \(K\), this gives the
polynomial dependence on \(N\) and \(e+1\) used in the rank-three
window construction, with explicit dependence here on the discriminant.
\src{src/Development/AutomorphicArithmetic/AutomorphicVarieties/UnitaryThreeGroups/KatoUse.lean}{52}{The fixed-field bound used on the dependency path}.
If two semisimple Galois representations are assigned the same full
Hecke system, equality of traces at almost all Frobenius elements,
Chebotarev and Brauer--Nesbitt make them isomorphic. For the reducible
list, assigning the same rank-two system gives isomorphic restrictions
and hence the same rank-three class. This supplies the injection on
each pairwise nonisomorphic sublist. For classical cuspidal attachment
and the stated local Galois conditions, Proposition~\ref{newgal:level}
transfers the principal level, including at the coefficient prime.
The outstanding requirement is a construction from every relevant strict
Selmer group which gives the length inequality and produces such
automorphic systems.
It is this construction which makes the integer relevant to cohomology.

\subsection{Quadratic restriction and the rank-four trace product}

In the rank-four problem the automorphy hypothesis over \(\mathbf Q\)
refers to continuous representations after restriction to imaginary
quadratic fields where the restriction is irreducible. Here is how a
bounded list over those fields can be used without choosing a
cuspidal representation over \(\mathbf Q\).
The following finite-list argument is an independent way to use this
part of the stronger automorphy hypothesis on the active path.
\src{src/Development/AutomorphicArithmetic/AutomorphicVarieties/ModularityArgument/ThresholdArguments/DenseThresholdCases/AutQ0.lean}{20}{Forgetting the additional automorphy witnesses};
\src{src/Development/AutomorphicArithmetic/AutomorphicVarieties/ModularityArgument/ThresholdArguments/DenseThresholdCases/VWLevel.lean}{18}{Use of level, automorphy and polarization in rank four}.

\begin{lemma}[Descent of the finite list]\label{rk4:quadratic}
Fix \(n\ge2\) and \(n^2+3\) distinct imaginary quadratic fields of
pairwise coprime discriminants. At a prime \(p\) discard the fields
ramified at \(p\); in each remaining field \(K_i\) choose a place
above \(p\), of residue degree \(f_i\in\{1,2\}\). Suppose finite lists
\(S_i\), independent of the coefficient prime, are given, with each
index \(s\) carrying an algebraic eigenvalue multiset \(B_{i,s}\)
of size \(n\), stable as a list under Galois conjugation.
For every \(\ell\ne p\), fix the coefficient embedding. Consider a
family of absolutely irreducible \(n\)-dimensional representations of
\(G_{\mathbf Q}\) unramified at \(p\). Whenever \(\ell\) is unramified
in \(K_i\) and the restriction is irreducible, suppose \(S_i\)
classifies the restrictions up to isomorphism and records their
arithmetic Frobenius eigenvalues at the chosen place. Then there is a
Galois-invariant list of at most \(2^{n+1}\sum_i|S_i|\) multiset
entries classifying the original representations with their arithmetic
Frobenius at \(p\). Here classification means an injective labelling
of isomorphism classes; different labels may have the same eigenvalues.
\end{lemma}
\begin{proof}
The prescribed fields can, for example, be obtained from distinct
primes \(q_i\equiv3\bmod4\) by adjoining \(\sqrt{-q_i}\).
For an irreducible representation \(r\), every quadratic self-twist
\(\eta\) gives a nonzero operator in \(\operatorname{End}(r)\) on
which conjugation acts by \(\eta\). Vectors belonging to distinct such
characters are independent. Indeed, apply \(g-\eta(g)\) to a shortest
nonzero relation, choosing \(g\) at which some other character differs
from \(\eta\); a shorter nonzero relation would result. Thus there are
at most \(n^2\) self-twists. For \(H=\ker\eta\) of index two, if
\(r|_H\) is reducible, take a simple \(H\)-submodule \(U\).
For \(g\notin H\), \(U+gU\) is stable under the whole group since
\(g^2\in H\). Irreducibility and properness of \(U\) imply it is
the direct sum \(U\oplus gU\). The operator acting by \(+1,-1\)
on the summands intertwines \(r\) with \(r\otimes\eta\).
At most one of the fields ramifies at \(p\), and at most one at
\(\ell\), so at least one of the remaining restrictions is irreducible.

If two extensions of an irreducible \(H\)-representation agree on \(H\)
in a chosen basis, the quotient of their matrices at \(g\notin H\)
commutes with \(r(H)\). Schur's lemma makes it scalar, and its square
is one since their matrices at \(g^2\) agree. Thus there are at most
two isomorphism classes of extensions. At the chosen place, arithmetic
Frobenius restricts as the \(f_i\)-th power of Frobenius at \(p\)
modulo inertia, up to conjugacy. Include for each \(i,s\) all multisets
\(A\) whose entrywise \(f_i\)-th powers form \(B_{i,s}\).
Ordering \(B_{i,s}\) temporarily shows there are at most
\(f_i^n\le2^n\) possible multisets, including when some entries
coincide. Give each two labels for the possible extensions. This is a
conjugation-stable finite list. Assign each representation to the first
eligible field in a fixed order, its restriction label, its actual
multiset and one of the two extension labels. This assignment is
injective for each fixed coefficient embedding. The whole list is
independent of that embedding and of \(\ell\).
\end{proof}

\begin{corollary}[The integer in rank four]\label{rk4:integer}
Take \(n=4\), \(p\ge5\), \(e\ge2\). Suppose, in the preceding lemma,
that every member of \(B_{i,s}\) is an algebraic integer of absolute
value \(p^{f_i(e+1)/2}\) in every complex embedding. Let \(S\) be
the resulting list and \(t_s\) its entrywise traces. Fix
\(z_1,z_2,z_3\in\mu_k\), \(k\ge1\). There is a nonzero rational
integer \(I_p\) with
\[
 \log|I_p|\le\varphi(k)|S|\bigl((e+1)\log p+\log 2\bigr)
\]
which bounds the sum of the \(v_\ell\)'s of the trace differences for
any pairwise nonisomorphic sublist at \(\ell\ne p\), for the comparison
trace \(p^{e+1}z_1+p^e z_2+pz_3+1\).
\end{corollary}
\begin{proof}
Every lifted eigenvalue is integral by the monic root equation,
and its absolute value is \(p^{(e+1)/2}\). For each
\(j\in(\mathbf Z/k)^\times\) put \(T_j=p^{e+1}z_1^j+p^ez_2^j+pz_3^j+1\);
at \(k=1\) use the unique automorphism of the trivial root-of-unity
group. Set \(P=p^{e+1}\). Since \((p^e+p+1)/P\le1/5+1/25+1/125<1/2\)
and \(P^{1/2}\ge5^{3/2}>8\), we have
\[
 |T_j|>P/2,\qquad |t_s|\le4\sqrt P<P/2,\qquad |t_s-T_j|<2P.
\]
The product \(I_p=\prod_{s,j}(t_s-T_j)\) is a nonzero algebraic
integer fixed by every Galois automorphism: the latter permutes the
three powers \emph{simultaneously} through the indicated indices.
This proves the size assertion. Every difference is integral at
\(\ell\); keeping just the factors belonging to a chosen sublist at
the identity automorphism gives the valuation bound.
\end{proof}
This gives the trace product in the arithmetic-Frobenius convention
of rank four. If the quadratic lists have size polynomial in \(M(e+1)\),
the bound has the required polynomial dependence, since only finitely
many fields were fixed for the rank. To apply it to a particular
Selmer group one needs the classical quadratic attachments of its
deformation points with the stated levels, weights and purity.
Section~\ref{spreal:section} identifies the branch on which an
attachment follows from a known modularity theorem. \src{src/Development/AutomorphicArithmetic/AutomorphicVarieties/ModularityArgument/ThresholdArguments/DenseThresholdCases/AfStmt.lean}{34}{Rank-four comparison required of those deformations}.
